\capitulo{6}{Quando um acerto já basta?}
Tenho duas turmas de aprendizes corrigindo os bilhetes da oficina. Quando as faces coincidem, todos escrevem A. Uma das turmas acha que A aparece em nove desses bilhetes a cada dez. A outra, em seis a cada dez. Se a prova pedir só a letra de cada bilhete, posso trocar as folhas de resposta das duas turmas e ninguém perceberá.

Você conhece um jeito de distingui-las. Volte à sala com o bilhete que custa R\$\,8,80 e paga dez reais se a máquina for A. Uma turma espera receber nove reais e compra. A outra espera seis e deixa passar. Suponha, porém, que a oficina só queira encaminhar os bilhetes a duas gavetas e que os dois tipos de erro custem o mesmo. Quanto vale, para esse serviço, distinguir duas apostas que abrem a mesma portinhola?

\secao{Apresente a conta por um acerto}
Vamos manter as faces, os rótulos e até a forma da função que usamos para classificar. Como no capítulo anterior, $z=+1$ significa faces iguais, e $z=-1$, faces diferentes. O escore será $f_w(z)=wz$. Identificamos A com $y=+1$ e B com $y=-1$. A simetria da nossa oficina permite usar o mesmo coeficiente nos dois tipos de bilhete. Para as faces iguais, dão A tanto um $w=0{,}4$ quanto um $w=4$. O que faz um desses números merecer o lugar do outro?

A regressão logística contém uma resposta. Seu escore mede as log-chances de A. Mas, ao corrigir a prova, já sabemos qual letra saiu. Um escore $-2$ quando sai B merece o mesmo crédito que um escore $+2$ quando sai A. Podemos virar a régua nos bilhetes B. Multiplicar o escore pelo rótulo $y$ deixa como está a leitura dos A e troca o sinal da leitura dos B. Chamemos de $m=yf_w(z)$ essa leitura do lado da letra que saiu. Nos dois acertos do exemplo, $m=2$.

Que probabilidade o aprendiz deu à letra certa? Na régua que acabamos de virar, as chances a favor dela são $e^m$ para um. Das $1+e^m$ partes, $e^m$ dizem respeito a ela. Sua fração é, portanto, $p_{\rm certa}=1/(1+e^{-m})$, como no capítulo anterior. Para encontrar a cobrança, procure no denominador quantas vezes teria de multiplicar essa fração para chegar a um. O logaritmo desse fator é o preço que instituímos no capítulo 2, agora em unidades naturais,
\[
             \ell_{\rm log}(m)=-\ln p_{\rm certa}=\ln(1+e^{-m}).
\]
Mesmo quem acertou a letra sai devendo alguma coisa. No grupo dos bilhetes iguais, o ponto mais barato ocorre quando anunciamos a fração de A que a pilha trouxe, como vimos no capítulo 4. Nove letras A para uma B pedem chances de nove para um, ou $w=\ln 9$. Se a pilha tiver seis A e quatro B, as chances pedidas são três para dois e o escore cai para $\ln(3/2)$. Sem penalização, a simetria das frações no grupo dos diferentes permite usar o mesmo coeficiente nos dois grupos.

Posso propor uma segunda correção. Vou cobrar um ponto pelo empate $m=0$. Se o aprendiz der um passo na direção certa, desconto um ponto. Se der um passo na direção errada, cobro um ponto a mais. A dívida tem um piso. De quem já chegou a $m=1$ não cobro mais nada naquele bilhete. Assim, um acerto com $m=0{,}6$ custa $0{,}4$ e um erro com $m=-0{,}6$ custa $1{,}6$. A regra é a perda \emph{hinge},
\[
                  \ell_{\rm h}(m)=\max(0,1-m).
\]
Chamamos o produto $m$ de \emph{margem funcional}. Vejamos o que a nova cobrança faz com aqueles nove acertos.

\perdasmargem

Ponha dez bilhetes com faces iguais sobre a mesa, nove de A e um de B. Com $w=0$, a conta da \emph{hinge} dá dez pontos. Ao subir de zero para um, você desconta um ponto em nove bilhetes e paga um ponto extra em um. Sobram dois pontos a pagar. Suponha que queira subir de um para dois. Qual dos nove bilhetes corretos ainda lhe dará desconto? Nenhum. A única conta que se mexe é a do erro, que encarece mais um ponto. A perda é linear em cada trecho, e andar para o lado negativo favoreceria a minoria. O melhor lugar para parar é $w=1$.

Troque três bilhetes de A por três de B. Restaram seis contra quatro. Até $w=1$, os seis descontos ainda pagam os quatro acréscimos. Depois de um, só os quatro erros encarecem. O melhor lugar continua sendo $w=1$. Posso lhe entregar o coeficiente vencedor, informar como escolhi seu valor e perguntar quanto devo pagar pelo próximo bilhete. Você teria como descobrir se eu estava com o maço dos nove ou com o dos seis? A portinhola lê o sinal positivo e dá a mesma resposta nos dois maços. Suponha agora que alguém aplique ao escore 1 a sigmoide da logística e anuncie cerca de $73\%$ para A. Essa operação daria o mesmo preço de R\$\,7,31 nos dois casos. Um maço sugere nove reais e o outro, seis. Posso até defender R\$\,7,31 se encontrar essa frequência em bilhetes novos com o mesmo escore. Mas essa promessa veio de uma hipótese a mais sobre o número 1. A perda \emph{hinge} só tinha cobrado a margem.

\secao{De onde vem a margem}
Para chegar à máquina de vetores de suporte, a \emph{SVM}, ainda falta fazer uma escolha. O inspetor da oficina traz um paquímetro. A portinhola agora separa peças pela medida do diâmetro. No lote de ajuste, a maior peça de B mede $4{,}8$ mm e a menor de A mede $5{,}2$ mm. Acima do limiar mandamos para A. Eu instalo o limiar em $5{,}15$ mm. Você instala o seu em $5{,}00$ mm. Testamos as peças uma a uma e os dois acertamos todas. Antes de mandar comprar a sua portinhola ou a minha, o que você gostaria de perguntar ao inspetor?

Ele pega a peça que mediu $5{,}2$, retira-a do paquímetro e torna a prendê-la. Conta que, só com a mudança de posição, duas leituras podem diferir até $0{,}10$ mm. Agora aparece uma diferença entre os dois acertos. Uma releitura de $5{,}10$ atravessa a minha portinhola para o lado de B. Na sua, a peça ainda tem folga. O mesmo vale do outro lado. Das fronteiras que acertaram o lote inteiro, a de $5{,}00$ mm é a que dá a maior folga à peça mais próxima de qualquer lado. A pessoa que sabia quanto valia essa folga estava segurando o paquímetro.

E por que não bastou premiar um escore grande do lado correto? Se $f(\boldsymbol u)=\boldsymbol w^\top\boldsymbol u+b$ é nossa régua, posso dobrar $\boldsymbol w$ e $b$ sem mover a fronteira $f=0$ nem um milímetro. Os escores dobraram. A distância da peça à fronteira continua onde estava. Para ligar as duas réguas, imagine que um milímetro perpendicular à fronteira valha quatro pontos no visor. Nesse caso, para ganhar um ponto basta andar um quarto de milímetro. A norma $\lVert\boldsymbol w\rVert$ faz o papel desses quatro pontos por milímetro. Multiplicar os deslocamentos nos eixos pelos pesos e somar mede o quanto caminhamos na direção de $\boldsymbol w$. Nessa direção, perpendicular à fronteira, uma distância $d$ acrescenta $d\lVert\boldsymbol w\rVert$ ao escore. Por isso, a linha de escore $+1$ fica a $1/\lVert\boldsymbol w\rVert$ da fronteira, e a linha de escore $-1$ fica à mesma distância do outro lado.

Quando conseguimos deixar todos os exemplos com $y_i f(\boldsymbol u_i)\geq1$, escolher a menor norma $\lVert\boldsymbol w\rVert$ abre a faixa mais larga. No caso do paquímetro, isso põe o limiar no meio das duas peças mais próximas. Já no bilhete com duas caras sabemos que há rótulos A e B para a mesmíssima entrada. Nenhuma reta os colocará em lados diferentes. Façamos então um acordo. Queremos abrir a faixa, mas aceitamos pagar pelas peças que caírem dentro dela ou do lado errado. Se eu dobrar o número de peças repetindo todas elas, quero que o preço médio dos erros continue igual. Por isso divido a cobrança das peças por $n$. Para o outro prato escolho cobrar pelo quadrado da norma. Dobrar a escala custa quatro vezes mais nessa parcela. É uma escolha sobre quanto vale a largura da faixa. Com essas duas cobranças, o custo que a SVM linear de \emph{margem suave} minimiza é
\[
 \frac1n\sum_{i=1}^n \max\{0,1-y_i f(\boldsymbol u_i)\}
       +\frac{\lambda}{2}\lVert\boldsymbol w\rVert^2,\qquad \lambda>0.
\]
O primeiro termo é a cobrança que fizemos nos dois maços. O segundo cobra por aumentar a escala dos coeficientes mesmo quando deixamos a fronteira exatamente onde estava. Um $\lambda$ maior dá mais peso à norma. A localização da fronteira depende das peças que alcançam a margem ou a violam, entre as quais ficam os chamados \emph{vetores de suporte}. O valor $w=1$ dos envelopes veio da perda sozinha. A penalização também participa do ajuste de uma SVM.\footnote{Corinna Cortes e Vladimir Vapnik desenvolvem o classificador com margem e violações em \href{https://doi.org/10.1007/BF00994018}{\emph{Support-Vector Networks}} (1995). Para a formulação por perda e penalização, veja T. Hastie, S. Rosset, R. Tibshirani e J. Zhu, \href{https://jmlr.org/papers/v5/hastie04a.html}{\emph{The Entire Regularization Path for the Support Vector Machine}} (2004).}

Na oficina eu testaria a SVM se quisesse uma separação estável sob releituras e conseguisse defender a régua em que medi a distância. Se eu der à máquina o diâmetro e o peso, trocar gramas por miligramas já muda a geometria dos dados crus. Ninguém aprendeu nada novo sobre as peças com essa troca de unidades. Cabe a mim escolher o que vale um deslocamento em cada eixo. Para verificar minha escolha, reservaria peças de origem conhecida e tornaria a medi-las. A folga que medi nas peças de ontem só me interessa se ajudar a encaminhar as que vão passar hoje.

Mude agora a encomenda. Queremos decidir se vale a pena inspecionar uma peça. A inspeção custa dois reais e evita um prejuízo de dez reais se ela tiver vindo de B. Chamando de $r$ o risco de B para aquela medida, vale inspecionar se $10r>2$. Sob um preço de quatro reais, o corte passa de $20\%$ para $40\%$. Uma peça com risco de $30\%$ muda de destino sem que ninguém mexa no paquímetro. Aqui eu começaria pela logística. Posso levar sua estimativa de risco à pessoa que sabe os preços. Há um detalhe. No lote que separou perfeitamente A e B, quem mediu o risco de uma peça de $5{,}05$ mm? Nenhuma peça tinha esse diâmetro. Uma curva logística cada vez mais íngreme pode melhorar o ajuste das letras do lote. Antes de usar seu número para gastar dinheiro, eu buscaria peças rotuladas perto da decisão e conferiria se os riscos anunciados se realizam. Também se pode aprender uma calibração para o escore de uma SVM. Para isso, eu separaria dados rotulados da população em que vou usar o escore e mediria ali a frequência de B.

\secao{Pode haver aprendizado quando a resposta continua igual?}
Deixe os dois programas em $w=2$ e entregue a ambos um bilhete repetido cuja letra é A. Eles acertaram. Agora deslize o dedo para a direita no gráfico das perdas, de $2$ para $2{,}01$. Na \emph{hinge} você andou sobre a parte horizontal. Na logística a curva ainda desce. Só de olhar você já sabe qual dos dois aprendizes recebe desse bilhete uma razão para aumentar $w$.

Quanto ele ainda tem a ganhar? Para $m=2$, a cobrança logística vale cerca de $0{,}1269$. Em $2{,}01$, vale $0{,}1257$. Um centésimo de passo economizou cerca de $0{,}0012$ ponto. A inclinação local deve andar perto de $-0{,}12$ ponto por unidade de margem. Podemos entender o tamanho desse número olhando a aposta. Na conta $\ln(1+e^{-m})$, o termo 1 fica parado. Se $m$ aumenta um centésimo, esse termo fica multiplicado por $e^{-0{,}01}$, cerca de $0{,}99$. Perde quase um centésimo de si mesmo. Para um passo cada vez menor, o tamanho que ele perde por unidade de passo tende ao próprio $e^{-m}$. O logaritmo mede essa mudança em relação ao total $1+e^{-m}$. A inclinação é, portanto,
\[
 \ell'_{\rm log}(m)=-\frac{e^{-m}}{1+e^{-m}}=-(1-p_{\rm certa}).
\]
O aprendiz que apostou perto de $88\%$ na letra certa ainda tinha quase $12\%$ do lado errado. É esse resto que o bilhete põe na conta do ajuste. Um \emph{passo de gradiente} manda andar contra a inclinação para baixar a despesa. Se $\eta>0$ regula o tamanho do passo, como aqui $m=w$, o bilhete acrescenta $\eta/(1+e^2)$ a $w$. Na \emph{hinge}, a inclinação zero acrescenta zero. Se cobrarmos também $\lambda w^2/2$ pelo coeficiente, por que um $w$ grande é puxado com mais força de volta a zero? Imagine $w^2$ como a área de um quadrado de lado $w>0$. Ao aumentar o lado um pouquinho, você acrescenta duas tiras de comprimento $w$ e um cantinho que encolhe com o quadrado do passo. A área cresce à taxa $2w$. Multiplique pela tarifa $\lambda/2$ e terá uma inclinação $\lambda w$. Essa parcela retira $\eta\lambda w$ dos dois passos. Assim conseguimos separar no movimento do botão quanto veio da nova letra e quanto veio da cobrança que já existia.

Poderíamos ter feito outra escolha desde o começo. Para cada tipo de bilhete, guarde apenas as contagens de A e B. Quando chegar uma letra nova, acrescente um à contagem correspondente. Com esse mesmo caderno posso responder qual letra é a mais frequente, calcular a proporção de A ou escolher o escore que minimiza cada uma das duas perdas naquela coluna. Mudar o que perguntei ao caderno custou uma nova conta na hora de lê-lo. A instrução de acrescentar um à contagem continuou servindo.

Na primeira página deste livro o estado do aprendiz era uma distribuição e a regra que a atualizava se chamava $\mathcal F$. O caderno cresceu. Em um momento ele guardou contagens e em outro, coeficientes de uma reta. Se $S_t$ é aquilo que guardamos, $d_t$ é o dado novo e $Q$ é a pergunta, temos lugar para duas operações,
\[
        S_t=U(S_{t-1};d_t),\qquad
        \text{resposta}=R_Q(S_t).
\]
Quando o estado é a distribuição sobre hipóteses, $U$ pode ser a nossa $\mathcal F$, e $R_Q$ pode simplesmente ler a probabilidade de A. A separação só se tornou necessária quando conseguimos fazer perguntas diferentes a um mesmo estado e dar o mesmo dado a regras que o tratam de maneiras diferentes.

\begin{quadro}{O que escolhi antes de abrir o caderno}
\itemq{Para apostar} Na logística escolho os atributos que compõem as log-chances e cobro a probabilidade da letra revelada. Para usar o resultado como risco, quero ver a frequência de erros nas peças em que desse risco depende uma decisão.
\itemq{Para separar} Na SVM escolho uma representação, uma régua de distância e o preço das violações da margem. Um exemplo classificado corretamente além da faixa já quitou sua parte na perda.
\end{quadro}

Ficou na mesa o maço com nove bilhetes de A e um de B. Você conhece só essas contagens e aposta pela frequência observada. Trago mais um bilhete de faces iguais. Viramos. É B. Você risca o dez no total da coluna e escreve onze. Pergunto pela próxima letra. Você diz A, como dizia antes. Ofereço então os dez reais por R\$\,8,80. Agora o bilhete lhe vale $10\times9/11$, cerca de R\$\,8,18. Você deixa a minha mão estendida.

Se eu tivesse anotado apenas a letra que você previu, teria passado a tarde inteira sem perceber que você aprendeu alguma coisa.
